\chapter{CONSIDERAÇÕES FINAIS}
\label{cap:conclusao}

O estágio supervisionado, correspondente à disciplina PRG -- 107,
realizado no setor de Clínica Médica de Equinos do Hospital Veterinário
da UFMG, possibilitou valorosa integração dos conhecimentos teóricos
adquiridos ao longo de toda a graduação com a vivência profissional
contemporânea. A observação e intensa inserção na rotina clínica do
setor permitiu o desenvolvimento e aperfeiçoamento de habilidades
técnicas e clínicas indispensáveis à formação do médico veterinário
especialista em equinos, bem como a responsabilidade ética e o
comprometimento necessários na condução dos referidos casos clínicos.

A possibilidade de interação com profissionais de diferentes formações e
com realidades clínicas diversas permitiu compreender que não existe um
tratamento universalmente ideal, mas sim aquele adequado à situação
clínica específica, considerando as condições do local, os recursos
disponíveis e sobretudo as particularidades de cada paciente e
proprietário. Aspectos como tomada de decisão e comunicação profissional
alinhada à prática veterinária foram práticas constantes durante esse
importante período de aprendizado na UFMG. O caso clínico relatado neste
trabalho, ainda que limitado pela extensão do estudo, ilustra com rigor
a complexidade dos desafios enfrentados na clínica médica de equinos,
reforçando assim a importância da formação técnica sólida, do raciocínio
clínico integrado, do senso crítico e da atualização constante frente às
questões emergentes de resistência antimicrobiana. A experiência
adquirida neste estágio foi de significativa importância para a
consolidação e aprimoramento da formação profissional pretendida e
reforçou consideravelmente o interesse da discente em seguir carreira na
área da medicina equina, associada ainda à pesquisa científica na
referida temática.
